%% 4-15 いろいろな座標：極座標・円柱座標・球座標（第13回：組 V4-7。ページ制限なし。偶数ページ）
%% ほかの巻から指されている形：3-2「原点からの距離と向きの角度の2つで表す方法もあります（極座標、4-15）」
%% 物理シリーズ（力学）のために必ず書く形：極座標と e_r・e_θ、円運動の速度 rω e_θ、だ円（軽く）
\makeatletter
\tikzset{ivp3d/.style={x={({-0.49*#1},{-0.34*#1})}, y={(#1,0)}, z={(0,#1)}}}
\newcommand\bkivpline[1]{\par\addvspace{3pt}{\centering #1\par}\addvspace{3pt}\noindent\ignorespaces}
\newcommand\bkivpstep[2]{\par\addvspace{6pt}\Needspace{5\baselineskip}\noindent
  {\bkcircled[bkOrange]{#1}\enspace\sffamily\bfseries\color{bkOrange!75!black}#2}\par\nobreak\vspace{2pt}}
\newcommand\bkivphead[1]{\par\addvspace{4pt}\Needspace{6\baselineskip}\noindent\textbf{#1}\par\nobreak}
\newcommand\bkivpaxes[3]{%
  \draw[bkGray, line width=0.6pt, -{Stealth[length=4pt]}] (0,0,0) -- (#1,0,0) node[anchor=east, inner sep=2pt, font=\sffamily\small, text=bkInk] {$x$};
  \draw[bkGray, line width=0.6pt, -{Stealth[length=4pt]}] (0,0,0) -- (0,#2,0) node[anchor=west, inner sep=2pt, font=\sffamily\small, text=bkInk] {$y$};
  \draw[bkGray, line width=0.6pt, -{Stealth[length=4pt]}] (0,0,0) -- (0,0,#3) node[anchor=south, inner sep=2pt, font=\sffamily\small, text=bkInk] {$z$};}
% 平面の座標軸（mm）  \bkivpplane{左}{右}{下}{上}
\newcommand\bkivpplane[4]{%
  \draw[bkGray, line width=0.6pt, -{Stealth[length=4pt]}] (#1,0) -- (#2,0) node[anchor=west, inner sep=2pt, font=\sffamily\small, text=bkInk] {$x$};
  \draw[bkGray, line width=0.6pt, -{Stealth[length=4pt]}] (0,#3) -- (0,#4) node[anchor=south, inner sep=2pt, font=\sffamily\small, text=bkInk] {$y$};}
% 地球儀（中心 #1、半径 #2 mm）
\newcommand\bkivpglobe[2]{\begin{scope}[shift={#1}]
  \fill[bkBlue!12] (0,0) circle (#2);
  \draw[bkBlue!70!black, line width=0.8pt] (0,0) circle (#2);
  \foreach \la in {-50,-25,25,50} {\draw[bkBlue!45, line width=0.4pt] ({-#2*cos(\la)},{#2*sin(\la)}) arc (180:360:{#2*cos(\la)} and {#2*cos(\la)*0.25});}
  \draw[bkRed!70!black, line width=0.7pt] (-#2,0) arc (180:360:#2 and {#2*0.25});
  \draw[bkRed!70!black, line width=0.4pt, dashed] (-#2,0) arc (180:0:#2 and {#2*0.25});
  \foreach \lo in {0.35,0.75} {\draw[bkBlue!45, line width=0.4pt] (0,-#2) arc (-90:90:{#2*\lo} and #2);}
  \draw[bkBlue!45, line width=0.4pt] (0,-#2) -- (0,#2);
\end{scope}}
\makeatother

\bkchapter{4}{15}{いろいろな座標：\\極座標・円柱座標・球座標}{\bklevelrange{4}{5}}{いろいろな座標：極座標・円柱座標・球座標}
\bkmichishirube{3-2,4-2,4-8,4-13}{8-6,8-12}

\begin{goal}
\begin{itemize}[label={\color{\bkTT}●}, leftmargin=1.6\zw, itemsep=0pt, topsep=1pt]
\item \bklevel{3}\enspace 極座標 $(r,\ \theta)$ と $(x,\ y)$ を、区画に注意して行き来できる。
\item \bklevel{4}\enspace 単位ベクトル $\vec e_r$、$\vec e_\theta$ を使って、円運動の速度を $r\omega\,\vec e_\theta$ と表せる。
\item \bklevel{5}\enspace 円柱座標・球座標で空間の点を表し、$(x,\ y,\ z)$ に直せる。
\end{itemize}
\end{goal}

\begin{hatena}
\noindent\begin{minipage}[c]{0.6\linewidth}
地図アプリで場所を調べると、「北緯35.7°、東経139.7°」のような2つの数が出てきます。これは東京のあたりの場所です。

地球の上の場所は、横へ何km・たてへ何kmではなく、「角度」で表しています。なぜ角度で表すのかな？　この2つの角度は、どこからはかった角度なのかな？
\end{minipage}\hfill
\begin{minipage}[c]{0.36\linewidth}
\centering
\begin{tikzpicture}[x=1mm, y=1mm]
  \bkivpglobe{(0,0)}{17}
  \fill[bkOrange] ({17*0.75*cos(52)},{17*sin(35.7)}) circle (1.3);
  \node[font=\sffamily\footnotesize, anchor=west, inner sep=2pt] at ({17*0.75*cos(52)+1},{17*sin(35.7)+1}) {東京};
  \node[font=\sffamily\footnotesize, text=bkRed!80!black, anchor=west] at (18,-1.5) {赤道};
  \node[font=\sffamily\footnotesize, anchor=south] at (0,17.5) {北極};
\end{tikzpicture}
\end{minipage}
\end{hatena}

\noindent 3-2では、平面の点を「右へいくつ、上へいくつ」の2つの数で表しました。4-2では、同じ場所を「どれだけはなれているか（距離）」と「どちらの向きか（角度）」で表す極座標を学び、4-6・4-8で2つの言い方を $\sin$・$\cos$ でつなぎました。この章では、その反対向きの直し方と、空間（3つの数）への広げ方を学びます。回るもの・丸いものは、距離と角度のほうがずっと楽に表せます。

%% ---------------- やってみよう・思い出そう ----------------
\begin{yattemiyou}[方眼の点の、距離と角度をはかろう]
\noindent\begin{minipage}[c]{0.56\linewidth}
\textbf{用意するもの}\par
方眼紙（1目もり\qty{1}{cm}）、分度器、定規
\begin{enumerate}[label={\bkcircled[bkGreen]{\arabic*}}, leftmargin=2.4\zw, itemsep=1pt, topsep=3pt]
\item 方眼紙のまん中に原点Oと、$x$ 軸・$y$ 軸をかこう。
\item 点A$(3,\ 4)$、B$(-3,\ 3)$、C$(-4,\ -3)$ をとろう。
\item 定規で、原点Oから、それぞれの点までの距離をはかろう。
\item 分度器で、$x$ 軸の＋の向きから左回りに、それぞれの点の向きの角をはかろう（180°をこえるときは、180°の向きからの角をはかって、180°にたす）。
\end{enumerate}
\end{minipage}\hfill
\begin{minipage}[c]{0.4\linewidth}
\centering
\begin{tikzpicture}[x=0.9mm, y=0.9mm]
  \draw[bkLine, line width=0.3pt] (-22,-22) grid[step=5] (22,22);
  \bkivpplane{-22}{24}{-22}{24}
  \foreach \px/\py/\n/\dx/\an in {15/20/A/1.6/west, -15/15/B/-1.6/east, -20/-15/C/-1.6/east} {
    \draw[bkRed, line width=0.9pt] (0,0) -- (\px,\py);
    \fill[bkRed] (\px,\py) circle (1);
    \node[font=\sffamily\footnotesize, anchor=\an, inner sep=1pt, fill=white] at ($(\px,\py)+(\dx,0)$) {\n};}
  \node[font=\sffamily\footnotesize, anchor=north east, inner sep=1pt, fill=white] at (-0.5,-0.5) {O};
\end{tikzpicture}
\end{minipage}
\par\vspace{10pt}\noindent ソラさんがはかった結果です（距離は\qty{0.1}{cm}まで、角は1°まで）。
\begin{bkcenter}
\begin{bktable}{colspec={Q[l,m,wd=12mm] Q[c,m,wd=17mm] Q[c,m,wd=17mm] Q[c,m,wd=19mm] Q[c,m,wd=19mm]}, rowsep=1.5pt, row{1}={font=\sffamily\small}}
点 & $x$ & $y$ & 距離 & 角度 \\
A & 3 & 4 & 5.0 & 53° \\
B & $-3$ & 3 & 4.2 & 135° \\
C & $-4$ & $-3$ & 5.0 & 217° \\
\end{bktable}
\end{bkcenter}
\end{yattemiyou}

\noindent 4-2で学んだように、同じ点は、$(x,\ y)$ でも、極座標 $(r,\ \theta)$ でも言えます。$(r,\ \theta)$ から $(x,\ y)$ へは、4-6・4-8の $x=r\cos\theta$、$y=r\sin\theta$ で計算できました。では反対に、$(x,\ y)$ から距離と角度を、はからずに計算で求めるには、どうすればよいでしょう。この章では、まずこの反対向きの直し方を作り、そのあと、空間の点を距離と角で表す方法に広げます。

\begin{omoidasou}
\bkrecall{3-2}{平面の点は $(x,\ y)$。\\原点と＋の向きは、問題を解く人が楽な所に決めてよい。\\符号で4つの区画（右上・左上・左下・右下）に分かれる。}
\bkrecall{4-2}{極座標 $(r,\ \theta)$：原点からの距離 $r$ と、$x$ 軸の＋の向きから左回りの角 $\theta$ で場所を言う。\\角は度でもラジアンでもはかる。$\pi$\,rad＝180°。弧の長さ＝半径×角（$s=r\theta$）。円の接線は半径に垂直。}
\bkrecall{4-8}{極座標 $(r,\ \theta)$ の点は、どの向きでも $x=r\cos\theta$、$y=r\sin\theta$（4-6で作り、4-8で広げた）。\\半径1の円の上の点は $(\cos\theta,\ \sin\theta)$。90°回すと $(x,\ y)\to(-y,\ x)$。$\theta=\omega t$、速さ $v=r\omega$。}
\bkrecall{4-13}{空間の点は $(x,\ y,\ z)$。原点からの距離は $\sqrt{x^2+y^2+z^2}$。}
\end{omoidasou}

\noindent この章では、$\approx$（「ほぼ等しい」というしるし。5-8でくわしく）も使います。

%% ---------------- 数学の道具 ----------------
\bksection{doguu}{数学の道具}

\bksubsection{1}{極座標と $(x,\ y)$ を行き来する}

\bkivpstep{1}{思い出そう：極座標と、$x=r\cos\theta$、$y=r\sin\theta$}\label{bk:4-15:kyoku}
\noindent\begin{minipage}[c]{0.5\linewidth}
\centering
\begin{tikzpicture}[x=1mm, y=1mm]
  \bkivpplane{-6}{50}{-6}{44}
  \draw[bkGray, line width=0.4pt, dashed] (30,0) -- (30,40) -- (0,40);
  \draw[bkRed, line width=1.6pt] (0,0) -- (30,40);
  \fill[bkInk] (30,40) circle (1.1);
  \node[font=\sffamily\small, anchor=south west, inner sep=1.5pt] at (30,40) {P};
  \draw[bkInk, line width=0.6pt, -{Stealth[length=3pt]}] (9,0) arc (0:53.13:9);
  \node[font=\sffamily\small] at (26.5:12.5) {$\theta$};
  \node[font=\sffamily\small, text=bkRed!80!black, anchor=south east, inner sep=1.5pt] at (14,21) {$r$};
  \node[font=\sffamily\footnotesize, anchor=north, inner sep=1.5pt] at (30,-0.5) {$x$};
  \node[font=\sffamily\footnotesize, anchor=east, inner sep=1.5pt] at (-0.5,40) {$y$};
  \node[font=\sffamily\footnotesize, anchor=north east, inner sep=1pt] at (-0.5,-0.5) {O};
  \draw[bkBlue, line width=1.1pt] (30,0) -- (30,40);
  \draw[bkGreen!70!black, line width=1.1pt] (0,0) -- (30,0);
\end{tikzpicture}
\end{minipage}\hfill
\begin{minipage}[c]{0.46\linewidth}
原点Oからの\textbf{距離} $r$ と、$x$ 軸の＋の向きから左回りにはかった\textbf{向きの角} $\theta$ の組 $(r,\ \theta)$ が、点Pの極座標でした（4-2）。原点Oを\textbf{極}、$x$ 軸の＋の部分を\textbf{始線}ともいいます。$r$ は0以上の数で、3-15の相関係数の $r$ とは別の意味です。
\end{minipage}

\par\vspace{4pt}
極座標から $(x,\ y)$ へは、斜辺 $r$ の直角三角形の横とたてで（4-6）、どの向きの点でも次の式で直せました（4-8）。
\bkivpline{\fbox{$x=r\cos\theta,\qquad y=r\sin\theta$}}
図のPは $r=5$、$\cos\theta=0.6$、$\sin\theta=0.8$ の点で、$x=5\times0.6=3$、$y=5\times0.8=4$ です。やってみようのAと同じ点です。この章では、この反対向き（$(x,\ y)$ から $(r,\ \theta)$）を考えます。

\bkivpstep{2}{同じ点に、いくつもの名前}
角は、1周（360°）回るともとの向きにもどります（4-8の一般角）。だから、極座標 $(4,\ 60°)$ と $(4,\ 420°)$ と $(4,\ -300°)$ は、同じ点です。ふつうは、$0°\leqq\theta<360°$（ラジアンなら $0\leqq\theta<2\pi$）の範囲で答えます。原点は $r=0$ で、角は決まりません（どの向きに0進んでも原点だから）。

\bkivpstep{3}{$(x,\ y)$ から極座標へ：区画に注意}
反対に、$(x,\ y)$ から $r$ と $\theta$ を求めます。$r$ は三平方の定理（4-4）で
\bkivpline{\fbox{$r=\sqrt{x^2+y^2}$}}
$\theta$ は、$\tan\theta=\dfrac{y}{x}$（傾き。4-6）なので、電卓の $\tan^{-1}$ で求められそうです。ところが、ここに落とし穴があります。

\noindent\begin{minipage}[c]{0.46\linewidth}
\centering
\begin{tikzpicture}[x=1mm, y=1mm]
  \bkivpplane{-30}{30}{-24}{30}
  \draw[bkLine, line width=0.3pt] (-25,-20) grid[step=6.67] (25,25);
  \draw[bkBlue, line width=1.4pt] (0,0) -- (-20,20);
  \fill[bkBlue] (-20,20) circle (1.2);
  \node[font=\sffamily\footnotesize, anchor=south, inner sep=2pt, fill=white] at (-20,21) {Q$(-3,\ 3)$};
  \draw[bkRed, line width=1.2pt, dashed] (0,0) -- (20,-20);
  \fill[bkRed] (20,-20) circle (1);
  \node[font=\sffamily\footnotesize, text=bkRed!80!black, anchor=north, inner sep=2pt, fill=white] at (20,-21) {$-45°$ の向き};
  \draw[bkBlue, line width=0.7pt, -{Stealth[length=3pt]}] (6,0) arc (0:135:6);
  \node[font=\sffamily\footnotesize, text=bkBlue!80!black, fill=white, inner sep=0.5pt] at (75:10.5) {135°};
\end{tikzpicture}
\end{minipage}\hfill
\begin{minipage}[c]{0.5\linewidth}
点Q$(-3,\ 3)$ で考えます。$\tan\theta=3\div(-3)=-1$ です。電卓で $\tan^{-1}(-1)$ を計算すると、$-45°$ と出ます。

\bkng{$\theta=-45°$}
ところが、$-45°$ の向きは右下（赤い点線）で、Qのある左上とは反対です。$\tan$ の値は、反対向き（180°ちがう向き）でも同じになってしまうのです。電卓の $\tan^{-1}$ は、いつも $-90°$ から $90°$ のあいだ（右半分）の角しか答えません。
\bkok{$\theta=-45°+180°=135°$}
\end{minipage}

\par\vspace{4pt}
だから、点がどの区画にあるかを見て、電卓の答えを直します。
\begin{bkcenter}
\begin{bktable}{colspec={Q[l,m,wd=24mm] Q[c,m,wd=28mm] Q[l,m,wd=56mm]}, rowsep=1.5pt, row{1}={font=\sffamily\small}}
点のある区画 & 符号 & $\theta$ の求め方（$0°\leqq\theta<360°$） \\
右上 & $x>0,\ y>0$ & $\tan^{-1}\dfrac{y}{x}$ のまま \\
左上・左下 & $x<0$ & $\tan^{-1}\dfrac{y}{x}$ に180°をたす \\
右下 & $x>0,\ y<0$ & $\tan^{-1}\dfrac{y}{x}$ に360°をたす \\
\end{bktable}
\bktabcaption{$x=0$ のとき（$y$ 軸の上）は、$y>0$ なら90°、$y<0$ なら270°。}
\end{bkcenter}
いちばん確かな方法は、\textbf{まず点を図にかいて、だいたいの向きを見てから}電卓を使うことです。表計算ソフトには、区画まで考えて角を出す関数 \texttt{ATAN2} があります。表計算では \texttt{ATAN2($x$, $y$)} の順に入れ、答えはラジアン（$-\pi$ より大きく $\pi$ 以下）で出ます（プログラムの言語では、$y$ を先に入れるものが多いので注意）。度にするには \texttt{DEGREES} 関数を使います（4-2）。

\bkivphead{もう1つの求め方：$\cos^{-1}$ と $y$ の符号}
$\cos\theta=\dfrac{x}{r}$ なので（$x=r\cos\theta$ の両辺を $r$ でわった）、電卓の $\cos^{-1}$ でも角が求められます。$\cos^{-1}$ は、いつも0°から180°（上半分）の角を答えます（4-6の⑧：$\cos$ の値から、0°〜180°の角は1つに決まる）。だから、
\begin{itemize}[leftmargin=1.2\zw, itemsep=1pt, topsep=2pt, label={\color{\bkTT}●}]
\item $y\geqq0$（上半分）なら、$\theta=\cos^{-1}\dfrac{x}{r}$ のまま。
\item $y<0$（下半分）なら、上下を折り返した向きなので、$\theta=360°-\cos^{-1}\dfrac{x}{r}$。
\end{itemize}
やってみようのC$(-4,\ -3)$ で、2つの方法をくらべます。$r=\sqrt{16+9}=5$ です。
\bkivpline{$\tan$ の方法：$\tan^{-1}\dfrac{-3}{-4}=\tan^{-1}0.75\approx36.87°$。左下なので $36.87°+180°=216.87°$}
\bkivpline{$\cos$ の方法：$\cos^{-1}\dfrac{-4}{5}=\cos^{-1}(-0.8)\approx143.13°$。$y<0$ なので $360°-143.13°=216.87°$}
どちらも約217°で、はかった角と合いました。どちらの方法でも、「点はどの区画にあるか」を見て、電卓の答えを直すことが大切です。

\bkivphead{\bkhatenaref の答え（その1）}
地球のような丸いものの上では、場所を「中心から見た向き」で表すと、とても楽です。緯度と経度は、地球の中心から見た向きの角度なのです。くわしくは、④の球座標で見ます。

\bksubsection{2}{動く向きの単位ベクトル：$\vec e_r$ と $\vec e_\theta$}

\bkivpstep{1}{場所ごとに向きが変わる単位ベクトル}\label{bk:4-15:er}
$(x,\ y)$ の座標では、$x$ の向きの単位ベクトル $\vec i$ と、$y$ の向きの $\vec j$ を使いました（4-13）。$\vec i$、$\vec j$ は、どの場所でも同じ向きです。

極座標では、点Pごとに、次の2つの向きを考えます。
\begin{itemize}[leftmargin=1.2\zw, itemsep=2pt, topsep=2pt, label={\color{\bkTT}●}]
\item $\vec e_r$：Pから、原点と反対の向き（$r$ がふえる向き）。長さ1。
\item $\vec e_\theta$：Pで、左回りに円にそう向き（$\theta$ がふえる向き）。長さ1。
\end{itemize}

\noindent\begin{minipage}[c]{0.5\linewidth}
\centering
\begin{tikzpicture}[x=1mm, y=1mm]
  \bkivpplane{-30}{32}{-8}{36}
  \draw[bkGray!70, line width=0.6pt] (0,0) circle (22);
  \foreach \a/\n in {35/P, 120/Q} {
    \draw[bkGray, line width=0.5pt] (0,0) -- (\a:22);
    \fill[bkInk] (\a:22) circle (1);
    \draw[bk arrow, bkRed, line width=1.5pt] (\a:22) -- ++(\a:10);
    \draw[bk arrow, bkBlue, line width=1.5pt] (\a:22) -- ++(\a+90:10);
    \node[font=\sffamily\small, anchor={\a+180}, inner sep=2pt] at (\a:22) {};}
  \node[font=\sffamily\footnotesize, anchor=north west, inner sep=1.5pt] at (35:22) {P};
  \node[font=\sffamily\footnotesize, anchor=north east, inner sep=1.5pt] at (120:22) {Q};
  \node[font=\sffamily\footnotesize, text=bkRed!80!black, anchor=west, inner sep=1pt] at ($(35:32)$) {$\vec e_r$};
  \node[font=\sffamily\footnotesize, text=bkBlue!80!black, anchor=south, inner sep=1.5pt] at ($(35:22)+(125:10)$) {$\vec e_\theta$};
  \node[font=\sffamily\footnotesize, text=bkRed!80!black, anchor=south, inner sep=1.5pt] at ($(120:32)$) {$\vec e_r$};
  \node[font=\sffamily\footnotesize, text=bkBlue!80!black, anchor=north east, inner sep=1pt] at ($(120:22)+(210:10)$) {$\vec e_\theta$};
\end{tikzpicture}
\end{minipage}\hfill
\begin{minipage}[c]{0.46\linewidth}
赤い矢印が $\vec e_r$、青い矢印が $\vec e_\theta$ です。PとQでは、矢印の向きがちがいます。\par\vspace{3pt}
$\vec e_r$ は、原点からPへの向きの単位ベクトルなので、4-8の単位円の点と同じで
\[ \vec e_r=(\cos\theta,\ \sin\theta) \]
$\vec e_\theta$ は、$\vec e_r$ を左回りに90°回したものです。
\end{minipage}

\par\vspace{4pt}
4-8で、点 $(x,\ y)$ を90°回すと $(-y,\ x)$ になりました。$\vec e_r=(\cos\theta,\ \sin\theta)$ を90°回すと、
\bkivpline{\fbox{$\vec e_r=(\cos\theta,\ \sin\theta),\qquad \vec e_\theta=(-\sin\theta,\ \cos\theta)$}}
たしかめます。$\vec e_r\cdot\vec e_\theta=\cos\theta\times(-\sin\theta)+\sin\theta\times\cos\theta=0$ なので、2つは直角です（4-12）。$|\vec e_\theta|^2=\sin^2\theta+\cos^2\theta=1$ なので、長さは1です（4-8）。$\theta=0°$ のPなら $\vec e_r=(1,\ 0)=\vec i$、$\vec e_\theta=(0,\ 1)=\vec j$ で、ふつうの $x$、$y$ の向きと同じになります。

点Pの位置ベクトル（原点からPへの矢印）は、長さ $r$ で $\vec e_r$ の向きなので、$\overrightarrow{\mathrm{OP}}=r\,\vec e_r$ と書けます。

\bkivpstep{2}{円運動の速度は $r\omega\,\vec e_\theta$}
半径 $r$ の円の上を、角速度 $\omega$（1秒に回る角、rad/s。4-8）で左回りに回る点Pを考えます。Pの速度は、
\begin{itemize}[leftmargin=1.2\zw, itemsep=2pt, topsep=2pt, label={\color{\bkTT}●}]
\item \textbf{速さ}\quad 1秒に回る角が $\omega$ なので、1秒に進む弧の長さは、弧の長さ＝半径×角（4-2）で $r\omega$。だから速さ $v=r\omega$（4-8）。
\item \textbf{向き}\quad 円にそって左回りに進むので、円の接線の向き（4-2で、接線は半径に垂直）。これは $\vec e_\theta$ の向き。
\end{itemize}
2つを合わせて、
\bkivpline{\fbox{円運動の速度\quad $\vec v=r\omega\,\vec e_\theta=(-r\omega\sin\theta,\ \ r\omega\cos\theta)$}}
$x$、$y$ の成分で書くと、$\vec v$ を $r\omega$ 倍した成分になります。

\bkivphead{数でたしかめる}
$r=\qty{0.5}{m}$、$\omega=\qty{4}{rad/s}$ で、いま $\theta=30°$ の所にいるとします。式では、速さ $v=0.5\times4=2$、$\vec e_\theta=(-\sin30°,\ \cos30°)=(-0.5,\ 0.866)$ なので、$\vec v=(-1,\ 1.732)$。

0.001秒だけ進めて、4-11の $\dfrac{\Delta x}{\Delta t}$、$\dfrac{\Delta y}{\Delta t}$ で調べます。0.001秒で角は $4\times0.001=0.004$\,rad ふえます。
\begin{bkcenter}
\begin{bktable}{colspec={Q[l,m,wd=34mm] Q[c,m,wd=30mm] Q[c,m,wd=30mm]}, rowsep=1.5pt}
 & $x=0.5\cos\theta$ & $y=0.5\sin\theta$ \\
はじめ（30°） & 0.433013 & 0.250000 \\
0.001秒後 & 0.432009 & 0.251730 \\
変化÷0.001 & $-1.00$ & $1.73$ \\
\end{bktable}
\bktabcaption{0.001秒のあいだの変化を0.001でわると、1秒あたりの変化（速度の成分）になる。}
\end{bkcenter}
式で求めた $(-1,\ 1.732)$ とよく合いました。$\vec e_\theta$ を使うと、円運動の速度を、図をかかずに1行で書けます。物理（力学）では、この形から、円運動の加速度を求めていきます。

\begin{motto}[4]
\textbf{速度の矢印の先も、円を回る}\par
円運動では、速さ $v=r\omega$ は変わりませんが、速度の向きは $\vec e_\theta$ といっしょに回ります。速度の矢印の根もとを1か所に集めると、矢印の先は、半径 $v$ の円の上を、同じ角速度 $\omega$ で回ります。だから、速度の先が動く速さは、「半径×角速度」（4-2・4-8）で $v\omega=r\omega^2$ です。速度が1秒あたりに変わる量（加速度）の大きさは $r\omega^2$ で、その向きは、速度の円の接線の向き、つまり $\vec e_\theta$ をさらに90°回した $-\vec e_r$（中心の向き）です。速さが一定でも、向きが変わるので、加速度があるのです（4-10の $\Delta\vec v\neq\vec0$）。$r=\qty{0.5}{m}$、$\omega=\qty{4}{rad/s}$ なら $0.5\times4^2=\qty{8}{m/s^2}$ です。
\end{motto}

\bksubsection{3}{円柱座標：極座標に高さを加える}

\label{bk:4-15:enchu}\noindent\begin{minipage}[c]{0.5\linewidth}
\centering
\begin{tikzpicture}[ivp3d=8mm]
  \bkivpaxes{3.4}{4.4}{5}
  % タワークレーン：柱は z 軸、ジブ（うで）は高さ 4.4
  \fill[bkYellow!60!black] (0,-0.12,0) -- (0,0.12,0) -- (0,0.12,4.6) -- (0,-0.12,4.6) -- cycle;
  \draw[bkYellow!50!black, line width=2.6pt] (0,0,4.4) -- ({2.8*cos(50)},{2.8*sin(50)},4.4);
  \draw[bkInk, line width=0.5pt] ({2.2*cos(50)},{2.2*sin(50)},4.4) -- ({2.2*cos(50)},{2.2*sin(50)},2);
  \fill[bkRed] ({2.2*cos(50)},{2.2*sin(50)},2) circle (1mm);
  \node[font=\sffamily\footnotesize, anchor=west, inner sep=2pt] at ({2.2*cos(50)},{2.2*sin(50)},2) {フック P};
  \draw[bkGray, line width=0.5pt, dashed] ({2.2*cos(50)},{2.2*sin(50)},2) -- ({2.2*cos(50)},{2.2*sin(50)},0);
  \draw[bkRed, line width=1.2pt] (0,0,0) -- ({2.2*cos(50)},{2.2*sin(50)},0);
  \draw[bkInk, line width=0.6pt, -{Stealth[length=3pt]}] plot[domain=0:50, samples=20] ({0.9*cos(\x)},{0.9*sin(\x)},0);
  \node[font=\sffamily\footnotesize, anchor=north, inner sep=2pt] at ({1.2*cos(25)},{1.2*sin(25)},0) {$\theta$};
  \node[font=\sffamily\footnotesize, text=bkRed!80!black, anchor=north west, inner sep=1.5pt] at ({1.4*cos(50)},{1.4*sin(50)},0) {$r$};
  \draw[bkGreen!70!black, line width=1.2pt] ({2.2*cos(50)},{2.2*sin(50)},0) -- ({2.2*cos(50)},{2.2*sin(50)},2);
  \node[font=\sffamily\footnotesize, text=bkGreen!50!black, anchor=east, inner sep=1.5pt] at ({2.2*cos(50)},{2.2*sin(50)},1) {$z$};
\end{tikzpicture}
\end{minipage}\hfill
\begin{minipage}[c]{0.46\linewidth}
ビルの工事現場のタワークレーンは、柱のまわりにうで（ジブ）を回し、うでにそってつり具を動かし、ワイヤーでフックを上げ下げします。フックの場所は、次の3つの数で決まります。
\begin{itemize}[leftmargin=1.2\zw, itemsep=1pt, topsep=2pt, label={\color{\bkTT}●}]
\item $r$：柱からの水平な距離
\item $\theta$：うでを回した角
\item $z$：高さ
\end{itemize}
\end{minipage}

\par\vspace{4pt}
床の上の位置を極座標 $(r,\ \theta)$ で表し、高さ $z$ を加えた組 $(r,\ \theta,\ z)$ を\textbf{円柱座標}（えんちゅうざひょう）といいます。\bkidx{えんちゅうざひょう}{円柱座標}$r$ を一定にして $\theta$ と $z$ を動かすと、点は円柱の面の上を動くので、この名前があります。$(x,\ y,\ z)$ には、①の式に高さをそのまま加えて直します。
\bkivpline{\fbox{$x=r\cos\theta,\qquad y=r\sin\theta,\qquad z=z$}}
円柱座標の $r$ は、原点からの距離ではなく、$z$ 軸からの距離です。原点からの距離は $\sqrt{r^2+z^2}$ になります（4-13の2段の三平方）。水道管・ドリル・モーターの軸のように、まっすぐな軸のまわりの形を表すときに便利です。

\bksubsection{4}{球座標：地球の上の場所}

\bkivpstep{1}{2つの角と、距離}\label{bk:4-15:kyu}
空間の点Pを、原点からの距離 $r$ と、2つの角で表します。
\begin{itemize}[leftmargin=1.2\zw, itemsep=2pt, topsep=2pt, label={\color{\bkTT}●}]
\item $r$：原点OからPまでの距離（4-13の $\sqrt{x^2+y^2+z^2}$）
\item $\theta$：$z$ 軸の＋の向き（真上）から、OPまでの角（$0°\leqq\theta\leqq180°$）
\item $\varphi$（ファイ）：Pの真下の床の点Qの、$x$ 軸から左回りの角（$0°\leqq\varphi<360°$）
\end{itemize}
この組 $(r,\ \theta,\ \varphi)$ を\textbf{球座標}（きゅうざひょう）といいます。\bkidx{きゅうざひょう}{球座標}$r$ を一定にすると、点は球の面の上を動きます。ここでの $\varphi$ は角の名前で、第5巻の初期位相の $\varphi$（5-15の $\phi$）とは別の意味です。

\noindent\begin{minipage}[c]{0.48\linewidth}
\centering
\begin{tikzpicture}[ivp3d=11mm]
  \bkivpaxes{3}{3.4}{3.2}
  \coordinate (P) at ({2.6*sin(50)*cos(40)},{2.6*sin(50)*sin(40)},{2.6*cos(50)});
  \coordinate (Q) at ({2.6*sin(50)*cos(40)},{2.6*sin(50)*sin(40)},0);
  \draw[bkRed, line width=1.6pt] (0,0,0) -- (P);
  \draw[bkGray, line width=0.6pt, dashed] (P) -- (Q);
  \draw[bkBlue, line width=1.2pt] (0,0,0) -- (Q);
  \fill[bkInk] (P) circle (0.9mm);
  \fill[bkInk] (Q) circle (0.7mm);
  \node[font=\sffamily\small, anchor=south west, inner sep=1.5pt] at (P) {P};
  \node[font=\sffamily\small, anchor=north west, inner sep=1.5pt] at (Q) {Q};
  \draw[bkInk, line width=0.6pt, -{Stealth[length=3pt]}] plot[domain=0:50, samples=20] ({0.9*sin(\x)*cos(40)},{0.9*sin(\x)*sin(40)},{0.9*cos(\x)});
  \node[font=\sffamily\small] at ({1.25*sin(22)*cos(40)},{1.25*sin(22)*sin(40)},{1.25*cos(22)}) {$\theta$};
  \draw[bkInk, line width=0.6pt, -{Stealth[length=3pt]}] plot[domain=0:40, samples=20] ({0.8*cos(\x)},{0.8*sin(\x)},0);
  \node[font=\sffamily\small, anchor=north, inner sep=2pt] at ({1.0*cos(20)},{1.0*sin(20)},0) {$\varphi$};
  \node[font=\sffamily\small, text=bkRed!80!black, anchor=north west, inner sep=1.5pt] at ($(0,0,0)!0.55!(P)$) {$r$};
\end{tikzpicture}
\end{minipage}\hfill
\begin{minipage}[c]{0.48\linewidth}
\textbf{$(x,\ y,\ z)$ に直す}\par
① Pの高さは、直角三角形OQPで、斜辺 $r$、真上とのなす角 $\theta$ なので
\[ z=r\cos\theta \]
② 床の上のOQ（青）の長さは $r\sin\theta$。\par\vspace{2pt}
③ Qは、床の上で、距離 $r\sin\theta$、角 $\varphi$ の点。①の極座標の式で
\[ x=r\sin\theta\cos\varphi,\quad y=r\sin\theta\sin\varphi \]
\end{minipage}

\par\vspace{4pt}
まとめると、
\bkivpline{\fbox{$x=r\sin\theta\cos\varphi,\qquad y=r\sin\theta\sin\varphi,\qquad z=r\cos\theta$}}
たしかめ：
\begin{align*}x^2+y^2+z^2&=r^2\sin^2\theta\,(\cos^2\varphi+\sin^2\varphi)+r^2\cos^2\theta\\&=r^2\sin^2\theta+r^2\cos^2\theta=r^2\end{align*}原点からの距離が $r$ になっています（$\sin^2+\cos^2=1$ を2回使った）。

反対に、$(x,\ y,\ z)$ からは、$r=\sqrt{x^2+y^2+z^2}$、$\cos\theta=\dfrac{z}{r}$ から $\theta$、Qの $(x,\ y)$ から①の方法で $\varphi$ を求めます。

\bkivpstep{2}{緯度・経度との関係}
地球の中心を原点にし、北極の向きを $z$ 軸、経度0°（イギリスのグリニッジを通る線）と赤道の交わる向きを $x$ 軸、東経90°の向きを $y$ 軸にとります（右手系）。

\noindent\begin{minipage}[c]{0.44\linewidth}
\centering
\begin{tikzpicture}[x=1mm, y=1mm]
  \fill[bkBlue!10] (0,0) circle (24);
  \draw[bkBlue!70!black, line width=0.8pt] (0,0) circle (24);
  \draw[bkRed!70!black, line width=0.7pt] (-24,0) -- (24,0);
  \node[font=\sffamily\footnotesize, text=bkRed!80!black, anchor=north east, inner sep=1pt] at (24,-0.5) {赤道};
  \draw[bkGray, line width=0.6pt, -{Stealth[length=3pt]}] (0,0) -- (0,31) node[anchor=south, font=\sffamily\footnotesize, text=bkInk] {$z$（北極）};
  \draw[bkInk, line width=1.2pt] (0,0) -- (35.7:24);
  \fill[bkOrange] (35.7:24) circle (1.2);
  \node[font=\sffamily\footnotesize, anchor=south west, inner sep=1.5pt] at (35.7:24) {東京};
  \draw[bkOrange!80!black, line width=0.7pt, -{Stealth[length=3pt]}] (9,0) arc (0:35.7:9);
  \node[font=\sffamily\footnotesize, text=bkOrange!80!black, anchor=west, inner sep=1pt] at (17.8:9.5) {緯度35.7°};
  \draw[bkGreen!70!black, line width=0.7pt, -{Stealth[length=3pt]}] (0,13) arc (90:35.7:13);
  \node[font=\sffamily\footnotesize, text=bkGreen!50!black, anchor=south west, inner sep=1pt] at (66:13) {$\theta$};
\end{tikzpicture}
\end{minipage}\hfill
\begin{minipage}[c]{0.52\linewidth}
\textbf{緯度}（いど）は、赤道から北（または南）へはかった角です。球座標の $\theta$ は真上（北極）からはかるので、北緯の角とは
\[ \theta=90°-\text{（北緯の角）} \]
の関係です。東京（北緯35.7°）なら $\theta=54.3°$。\par\vspace{3pt}
\textbf{経度}（けいど）は、経度0°の線から東へはかった角で、球座標の $\varphi$ と同じです。東経139.7°なら $\varphi=139.7°$。
\end{minipage}

\par\vspace{4pt}
$\sin(90°-\beta)=\cos\beta$、$\cos(90°-\beta)=\sin\beta$（4-6）を使うと、北緯 $\beta$、東経 $\varphi$ の地面の点は、地球の半径を $R$ として
\bkivpline{\fbox{$x=R\cos\beta\cos\varphi,\qquad y=R\cos\beta\sin\varphi,\qquad z=R\sin\beta$}}
と書けます。$R\cos\beta$ は、北緯 $\beta$ の緯線（同じ緯度の点を結んだ円）の半径です。

\bkivphead{\bkhatenaref の答え（その2）}
地球はほぼ球なので、中心からの距離 $r$ は、地面の上ならどこでも同じ地球の半径（約\qty{6400}{km}）です。だから、場所を決めるのに必要なのは、残りの2つの角だけです。1つは赤道からはかった北や南への角（緯度）、もう1つは経度0°の線からはかった東や西への角（経度）です。球座標で、$r$ がいつも同じだから、2つの角度だけで場所が言えるのです。

\bksubsection{5}{曲線を、極座標の式で表す}

\leavevmode\label{bk:4-15:kyokuhou}$x$ と $y$ の関係を式で書くと、グラフ（曲線）が決まりました（3-3）。同じように、$r$ と $\theta$ の関係の式でも、曲線が決まります。これを\textbf{極方程式}といいます。

\begin{bkfigure}
\begin{tikzpicture}
\begin{axis}[bkplot, name=g1, width=44mm, height=44mm, axis equal image, xmin=-4, xmax=4.4, ymin=-4, ymax=4.4, xtick=\empty, ytick=\empty, title={\sffamily\small 円 $r=3$}, title style={yshift=2mm}]
  \addplot[bkcurve, domain=0:360, samples=91] ({3*cos(x)},{3*sin(x)});
  \node[font=\sffamily\footnotesize, anchor=north west, inner sep=1.5pt] at (axis cs:3,0) {3};
\end{axis}
\begin{axis}[bkplot, name=g2, at={($(g1.east)+(8mm,0)$)}, anchor=west, width=44mm, height=44mm, axis equal image, xmin=-5.5, xmax=6.8, ymin=-5.5, ymax=6.8, xtick=\empty, ytick=\empty, title={\sffamily\small らせん $r=0.5\theta$}, title style={yshift=2mm}]
  \addplot[bkcurve2, domain=0:720, samples=181] ({0.5*x*pi/180*cos(x)},{0.5*x*pi/180*sin(x)});
\end{axis}
\begin{axis}[bkplot, name=g3, at={($(g2.east)+(8mm,0)$)}, anchor=west, width=44mm, height=44mm, axis equal image, xmin=-7, xmax=3.4, ymin=-4.8, ymax=4.8, xtick=\empty, ytick=\empty, title={\sffamily\small だ円 $r=\dfrac{3}{1+0.5\cos\theta}$}, title style={yshift=1mm}]
  \addplot[bkcurve3, domain=0:360, samples=181] ({3/(1+0.5*cos(x))*cos(x)},{3/(1+0.5*cos(x))*sin(x)});
  \begin{pgfonlayer}{axis foreground}
  \fill[bkRed] (axis cs:0,0) circle (2.5pt);
  \node[font=\sffamily\footnotesize, text=bkRed!80!black, anchor=south east, inner sep=1pt, fill=white] at (axis cs:-0.35,0.35) {焦点};
  \end{pgfonlayer}
  \node[font=\sffamily\footnotesize, anchor=north west, inner sep=1.5pt] at (axis cs:2,0) {2};
  \node[font=\sffamily\footnotesize, anchor=north east, inner sep=1.5pt] at (axis cs:-6,0) {$-6$};
\end{axis}
\end{tikzpicture}
\bkfigcaption{左：$r$ が一定なら円。中：角がふえるほど $r$ がふえると、らせん（$\theta$ はラジアン、2周分）。右：原点（赤い点）が、だ円の「焦点」になる。}
\end{bkfigure}

\begin{itemize}[leftmargin=1.2\zw, itemsep=3pt, topsep=2pt, label={\color{\bkTT}●}]
\item \textbf{円 $r=3$}\quad 原点からの距離がいつも3の点の集まりなので、半径3の円です。$(x,\ y)$ で書くと $x^2+y^2=9$（4-4）。極座標のほうが、ずっと短く書けます。
\item \textbf{らせん $r=0.5\theta$}\quad 1回りする（$\theta$ が $2\pi$ ふえる）ごとに、$r$ が $0.5\times2\pi=\pi$（約3.14）ずつふえる、等間かくのうずまきです。かとり線香の形や、時計のぜんまい、昔のレコードの溝は、このような形です。
\item \textbf{だ円 $r=\dfrac{3}{1+0.5\cos\theta}$}\quad $\theta=0°$ で $r=\dfrac{3}{1.5}=2$、$\theta=90°$ で $r=3$、$\theta=180°$ で $r=\dfrac{3}{0.5}=6$。原点から見て、右には近く、左には遠い、つぶれた円（だ円）になります。惑星は、太陽をこのような点（焦点）においた、だ円の上を回っています（ケプラーの法則。物理シリーズで学びます）。
\end{itemize}
4-8では、だ円を $(a\cos t,\ b\sin t)$ と表しました。こちらは、だ円の中心が原点にある表し方です。図のだ円（3つの図の右）のように、焦点を原点にすると、極方程式で表せます。

\begin{motto}[4]
\textbf{表計算で、極方程式のグラフをかこう}\par
極方程式 $r=1+\cos\theta$ のグラフを、表計算でかきます。表計算の \texttt{COS}・\texttt{SIN} はラジアンで計算するので、\texttt{RADIANS} で直してから入れます（4-2）。2行めに次の式を入れ、下へコピーします。
\begin{itemize}[leftmargin=1.2\zw, itemsep=0pt, topsep=2pt, label={\color{\bkTT}●}]
\item B2（$r$）：\texttt{=1+COS(RADIANS(A2))}
\item C2（$x$）：\texttt{=B2*COS(RADIANS(A2))}\quad D2（$y$）：\texttt{=B2*SIN(RADIANS(A2))}
\end{itemize}
\par\vspace{2pt}\noindent\begin{minipage}[c]{0.56\linewidth}
\centering
\bksheet{4}{11}{%
1 & $\theta$ & $r$ & $x$ & $y$ \\
2 & 0 & 2 & 2 & 0 \\
3 & 60 & 1.5 & 0.75 & 1.299 \\
4 & 120 & 0.5 & $-0.25$ & 0.433 \\
5 & 180 & 0 & 0 & 0 \\
6 & 240 & 0.5 & $-0.25$ & $-0.433$ \\
7 & 300 & 1.5 & 0.75 & $-1.299$ \\}
\end{minipage}\hfill
\begin{minipage}[c]{0.4\linewidth}
\centering
\begin{tikzpicture}
\begin{axis}[bkplot, width=46mm, height=42mm, axis equal image, xmin=-0.8, xmax=2.7, ymin=-1.6, ymax=1.7, xtick={1,2}, ytick={-1,1}, xticklabels={}, yticklabels={}]
  \addplot[bkcurve, domain=0:360, samples=145] ({(1+cos(x))*cos(x)},{(1+cos(x))*sin(x)});
  \addplot[only marks, mark=*, mark size=1.6pt, bkOrange] coordinates {(2,0) (0.75,1.299) (-0.25,0.433) (0,0) (-0.25,-0.433) (0.75,-1.299)};
  \node[font=\sffamily\footnotesize, anchor=north, inner sep=1.5pt] at (axis cs:1,-0.05) {1};
  \node[font=\sffamily\footnotesize, anchor=north west, inner sep=1pt] at (axis cs:2.2,-0.15) {2};
\end{axis}
\end{tikzpicture}
\end{minipage}\par\vspace{2pt}
A列の角（度）を5°ごとにして、C列とD列で散布図をかくと、図のような形（オレンジの点が表の6点）になります。この曲線をカージオイド（心臓の形）といいます。$\theta=180°$ で $r=0$ になり、原点を通ります。
\end{motto}

\bksubsection{6}{どの座標を使うか：楽になるものを選ぶ}

\leavevmode\label{bk:4-15:erabu}3-2では、「原点と＋の向きは、問題を解く人が楽な所に決めてよい」と学びました。座標の種類も、同じように自分で選べます。同じ場所を、どの座標で言っても、答えは同じです。計算が楽になるものを選ぶのが、上手な使い方です。
\begin{bkcenter}
\begin{bktable}{colspec={Q[l,m,wd=30mm] Q[l,m,wd=46mm] Q[l,m,wd=42mm]}, rowsep=2pt, row{1}={font=\sffamily\small}}
座標 & 楽になるもの & 例 \\
$(x,\ y,\ z)$ & 箱の形、まっすぐな動き & 部屋、ビル、3Dプリンター \\
極座標 $(r,\ \theta)$ & 1点のまわりを回るもの & レーダー、回転、円運動 \\
円柱座標 $(r,\ \theta,\ z)$ & まっすぐな軸のまわりの形 & 水道管、クレーン、モーター \\
球座標 $(r,\ \theta,\ \varphi)$ & 1点から四方に広がるもの & 地球、電波、原子の中の電子 \\
\end{bktable}
\bktabcaption{形に合わせて座標を選ぶ。}
\end{bkcenter}
たとえば、スマートフォンの電波は、アンテナから四方へ広がります。電波の強さは、アンテナからの距離 $r$ だけで決まり、向きにはあまりよりません（理想的な場合）。球座標なら「$r$ だけの式」で書けますが、$(x,\ y,\ z)$ では $\sqrt{x^2+y^2+z^2}$ という長い式になります。第8巻（8-6・8-12）では、円柱座標や球座標を使って、パイプの中の流れや、球のまわりの電場を計算します。

\begin{senmonsho}[極座標の書き方]
\begin{bkhonbun}
平面極座標 $(r,\theta)$ において、単位ベクトルを
\[ \symbf e_r=(\cos\theta,\sin\theta),\qquad \symbf e_\theta=(-\sin\theta,\cos\theta) \]
とすると、位置ベクトルは $\symbf r=r\symbf e_r$ である。半径 $r$ の等速円運動では、速度は $\symbf v=r\omega\,\symbf e_\theta$、加速度は $-r\omega^2\symbf e_r$ となる。球座標 $(r,\theta,\phi)$ との関係は
\[ x=r\sin\theta\cos\phi,\quad y=r\sin\theta\sin\phi,\quad z=r\cos\theta \]
\end{bkhonbun}
\bkyaku{太字はベクトル（4-13）。$\phi$ は $\varphi$ と同じ文字（ファイ）の別の形。数学の本の中には、$\theta$ と $\varphi$ の役わりを入れかえて書くものもあるので、どちらの角がどこからはかった角かを、まず確かめる。}
\end{senmonsho}

\begin{column}{飛行機は、なぜ北へふくらんで飛ぶ？}
東京からアメリカのニューヨークへ行く飛行機の道すじを、ふつうの世界地図にかくと、北へ大きくふくらんだ曲線になります。遠回りに見えますが、じつはこれが、地球の表面にそった最短の道すじです。球の上の2点を結ぶ最短の道は、地球の中心を通る平面で切った大きな円（大円）の一部になります。平らな地図は、丸い地球を引きのばしてかいているので、最短の道が曲がって見えるのです。例題7で、この距離を計算してみましょう。
\end{column}

\begin{taisetsu}
\def\tsl#1{\makebox[5.6\zw][l]{\sffamily\bfseries\small\color{\bkTT}#1}}
\tsl{極座標}$x=r\cos\theta$、$y=r\sin\theta$\\[3pt]
\tsl{}$r=\sqrt{x^2+y^2}$。$\theta$ は区画を見て決める\\[3pt]
\tsl{単位ベクトル}$\vec e_r=(\cos\theta,\ \sin\theta)$、$\vec e_\theta=(-\sin\theta,\ \cos\theta)$\\[3pt]
\tsl{円運動}速度 $\vec v=r\omega\,\vec e_\theta$（速さ $r\omega$、向きは接線）\\[3pt]
\tsl{円柱座標}$(r,\ \theta,\ z)$：$x=r\cos\theta$、$y=r\sin\theta$、$z=z$\\[3pt]
\tsl{球座標}$(r,\ \theta,\ \varphi)$：$x=r\sin\theta\cos\varphi$、$y=r\sin\theta\sin\varphi$、$z=r\cos\theta$\\[3pt]
\tsl{緯度と経度}$\theta=90°-$（北緯）、$\varphi=$（東経）
\end{taisetsu}

\begin{kotoba}
\small
\bkword{極座標}{原点からの距離 $r$ と、$x$ 軸からの角 $\theta$ で点を表す座標。}
\bkword{円柱座標}{極座標に、高さ $z$ を加えた座標 $(r,\ \theta,\ z)$。}
\bkword{球座標}{原点からの距離 $r$、真上からの角 $\theta$、床の上の角 $\varphi$ で表す座標。}
\bkword{緯度・経度}{地球の上の場所を表す2つの角。緯度は赤道から、経度は経度0°の線からはかる。}
\bkword{極方程式}{$r$ と $\theta$ の関係で曲線を表した式。}
\end{kotoba}

\bkkigocards{\bkkigocard{$(r,\ \theta)$}{アール・シータ}{極座標。原点からの距離と向きの角。$r$ はほかの章では別の意味（相関係数など）にも使う}{polar coordinates}{4-2}}
{\bkkigocard{\scalebox{0.86}{$(r,\ \theta,\ \varphi)$}}{アール・シータ・ファイ}{球座標。$\varphi$ は第5巻では初期位相の意味にも使う}{spherical coordinates}{4-15}}

%% ---------------- 例題 ----------------
\bksection{reidai}{例題でたしかめよう}

\begin{reidai}[$(x,\ y)$ から極座標へ（4つの区画）]
次の点の極座標 $(r,\ \theta)$ を求めましょう（$0°\leqq\theta<360°$）。\par
（1）$(1,\ \sqrt3)$\quad（2）$(-3,\ 3)$\quad（3）$(-2,\ -2\sqrt3)$\quad（4）$(2,\ -2)$\par
\kangaekata $r=\sqrt{x^2+y^2}$。$\theta$ は、まず区画を見てから、$\tan^{-1}\dfrac{y}{x}$ を直す（①の表）。\par
\shiki（1）$r=\sqrt{1+3}=2$。右上。$\tan^{-1}\sqrt3=60°$ のまま\par
\hspace*{2.6\zw}（2）$r=\sqrt{9+9}=3\sqrt2$。左上。$\tan^{-1}(-1)=-45°$\par
\hspace*{4.6\zw}$-45°+180°=135°$\par
\hspace*{2.6\zw}（3）$r=\sqrt{4+12}=4$。左下。$\tan^{-1}\sqrt3=60°$、$60°+180°=240°$\par
\hspace*{2.6\zw}（4）$r=\sqrt{4+4}=2\sqrt2$。右下。$\tan^{-1}(-1)=-45°$\par
\hspace*{4.6\zw}$-45°+360°=315°$\par
\kotae （1）$(2,\ 60°)$\quad（2）$(3\sqrt2,\ 135°)$\quad（3）$(4,\ 240°)$\quad（4）$(2\sqrt2,\ 315°)$
\end{reidai}
（3）では、$\dfrac{y}{x}=\dfrac{-2\sqrt3}{-2}=\sqrt3$ で、右上の（1）と同じ $\tan$ の値になります。区画を見ないと、まちがえてしまいます。

\begin{jibun}
次の点の極座標を求めましょう（$0°\leqq\theta<360°$）。\par
（1）$(\sqrt3,\ 1)$\quad（2）$(-1,\ 1)$\quad（3）$(-3,\ -3\sqrt3)$\quad（4）$(3,\ -3)$
\bkjans{（1）\nobreak $(2,\ 30°)$。（2）\nobreak $(\sqrt2,\ 135°)$。（3）\nobreak $(6,\ 240°)$。（4）\nobreak $(3\sqrt2,\ 315°)$}
\end{jibun}

\begin{tsumazuki}{$\tan^{-1}$ の答えを、そのまま角にする}
\bkng{$(-3,\ 3)$ の角は $\tan^{-1}(-1)=-45°$}
\bkok{左上の区画なので $-45°+180°=135°$}
$\tan$ は、180°ちがう向きでも同じ値になります。点を図にかいて、区画をたしかめてから角を決めましょう。
\end{tsumazuki}

\begin{reidai}[円運動の速度を $\vec e_\theta$ で]
半径\qty{0.5}{m}の円の上を、角速度\qty{4}{rad/s}で左回りに回る点が、いま $\theta=30°$ の所にいます。点の位置 $(x,\ y)$ と、速度の成分を求めましょう。\par
\kangaekata 位置は $(r\cos\theta,\ r\sin\theta)$。速度は $r\omega\,\vec e_\theta$、$\vec e_\theta=(-\sin\theta,\ \cos\theta)$（②）。\par
\shiki 位置\enspace $x=0.5\cos30°=0.5\times0.866=0.433$\par
\hspace*{5.6\zw}$y=0.5\sin30°=0.5\times0.5=0.25$\par
\hspace*{2.6\zw}速さ\enspace $r\omega=0.5\times4=2$\par
\hspace*{2.6\zw}$\vec e_\theta=(-\sin30°,\ \cos30°)=(-0.5,\ 0.866)$\par
\hspace*{2.6\zw}$\vec v=2\times(-0.5,\ 0.866)=(-1,\ 1.732)$\par
\kotae 位置 $(0.433,\ 0.25)$（m）、速度 $(-1,\ 1.73)$（m/s）
\end{reidai}

\begin{jibun}
半径\qty{2}{m}の円の上を、角速度\qty{3}{rad/s}で左回りに回る点が、いま $\theta=120°$ の所にいます。位置と速度の成分を求めましょう。
\bkjans{位置 $(-1,\ 1.73)$。速さ6、$\vec e_\theta=(-0.866,\ -0.5)$。速度 $(-5.20,\ -3)$（m/s）}
\end{jibun}

\begin{reidai}[円柱座標]
（1）タワークレーンのフックが、円柱座標で $r=\qty{20}{m}$、$\theta=60°$、$z=\qty{30}{m}$ の所にあります。$(x,\ y,\ z)$ を求めましょう。\par
（2）点 $(-6,\ 8,\ 5)$ の円柱座標を求めましょう（$\theta$ は電卓で）。\par
\kangaekata $x=r\cos\theta$、$y=r\sin\theta$、$z=z$（③）。反対は、床の上の $(x,\ y)$ を極座標にする（①）。\par
\shiki（1）$x=20\cos60°=10$、\enspace $y=20\sin60°=20\times0.866=17.32$\par
\hspace*{4.6\zw}$z=30$（高さはそのまま）\par
\hspace*{2.6\zw}（2）$r=\sqrt{36+64}=10$。左上なので\par
\hspace*{4.6\zw}$\tan^{-1}\dfrac{8}{-6}\approx-53.13°$、\enspace $-53.13°+180°=126.87°$\par
\kotae （1）約 $(10,\ 17.3,\ 30)$（m）\quad（2）$(10,\ 126.9°,\ 5)$
\end{reidai}

\begin{jibun}
（1）円柱座標 $r=10$、$\theta=150°$、$z=20$ の点の $(x,\ y,\ z)$ を求めましょう。\par
（2）点 $(5,\ -5,\ 12)$ の円柱座標を求めましょう。
\bkjans{（1）\nobreak 約 $(-8.66,\ 5,\ 20)$。（2）\nobreak $r=5\sqrt2\approx7.07$、右下で $-45°+360°=315°$。$(5\sqrt2,\ 315°,\ 12)$}
\end{jibun}

\begin{tsumazuki}{円柱座標の $r$ を、原点からの距離だと思う}
\bkng{円柱座標 $(3,\ 0°,\ 4)$ の点の原点からの距離は3}
\bkok{$r=3$ は $z$ 軸からの距離。原点からは $\sqrt{3^2+4^2}=5$}
円柱座標の $r$ は床の上の距離（$z$ 軸からの距離）です。球座標の $r$ とは意味がちがいます。
\end{tsumazuki}

\begin{reidai}[球座標]
（1）球座標 $r=2$、$\theta=60°$、$\varphi=45°$ の点の $(x,\ y,\ z)$ を求めましょう。\par
（2）点 $(1,\ 1,\ \sqrt2)$ の球座標を求めましょう。\par
\kangaekata $x=r\sin\theta\cos\varphi$、$y=r\sin\theta\sin\varphi$、$z=r\cos\theta$（④）。反対は、$r$、$\cos\theta=\dfrac{z}{r}$、床の上の点の角 $\varphi$ の順。\par
\shiki（1）$r\sin\theta=2\sin60°=2\times0.866=1.732$\par
\hspace*{4.6\zw}$x=1.732\cos45°=1.732\times0.7071\approx1.22$、\enspace $y\approx1.22$\par
\hspace*{4.6\zw}$z=2\cos60°=1$\par
\hspace*{2.6\zw}（2）$r=\sqrt{1+1+2}=2$。\enspace $\cos\theta=\dfrac{\sqrt2}{2}$ なので $\theta=45°$\par
\hspace*{4.6\zw}床の上の点 $(1,\ 1)$ は右上で、$\varphi=\tan^{-1}1=45°$\par
\kotae （1）約 $(1.22,\ 1.22,\ 1)$\quad（2）$(2,\ 45°,\ 45°)$
\end{reidai}
たしかめ：（1）の点の原点からの距離は $\sqrt{1.22^2+1.22^2+1^2}=\sqrt{3.98}\approx2.0$ で、$r=2$ と合います。答えが出たら、$r$ にもどるかを見るとまちがいに気づけます。

\begin{jibun}
（1）球座標 $r=4$、$\theta=30°$、$\varphi=120°$ の点の $(x,\ y,\ z)$ を求めましょう。\par
（2）点 $(0,\ 2,\ 2\sqrt3)$ の球座標を求めましょう。
\bkjans{（1）\nobreak $r\sin\theta=2$。約 $(-1,\ 1.73,\ 3.46)$。（2）\nobreak $r=4$、$\cos\theta=\dfrac{\sqrt3}{2}$ で $\theta=30°$、床の上の点 $(0,\ 2)$ は $y$ 軸の上で $\varphi=90°$。$(4,\ 30°,\ 90°)$}
\end{jibun}

\begin{reidai}[緯度・経度と地球の自転]
地球を半径\qty{6400}{km}の球とします。\par
（1）北緯30°、東経60°の地点の $(x,\ y,\ z)$ を求めましょう。\par
（2）北緯35.7°の東京で、地球の自転によって地面が東へ動く速さは、時速何kmですか。地球は24時間で1回転するとします。\par
\kangaekata（1）$x=R\cos\beta\cos\varphi$、$y=R\cos\beta\sin\varphi$、$z=R\sin\beta$（④）。（2）東京は、半径 $R\cos35.7°$ の緯線の円の上を、24時間で1周する。\par
\shiki（1）$R\cos30°=6400\times0.866=5542.6$\par
\hspace*{4.6\zw}$x=5542.6\times\cos60°\approx2771$、\enspace $y=5542.6\times\sin60°\approx4800$\par
\hspace*{4.6\zw}$z=6400\times\sin30°=3200$\par
\hspace*{2.6\zw}（2）緯線の半径\enspace $6400\times\cos35.7°\approx6400\times0.812=5197$\par
\hspace*{4.6\zw}1周\enspace $2\pi\times5197\approx\num{32655}$（km）\par
\hspace*{4.6\zw}$\qty{32655}{km}\div\qty{24}{h}\approx\qty{1361}{km/h}$\par
\kotae （1）約 $(2771,\ 4800,\ 3200)$（km）\quad（2）時速約\qty{1360}{km}
\end{reidai}
秒速に直すと、$\qty{1361}{km/h}\times\qty{1000}{m/km}\div\qty{3600}{s/h}\approx\qty{378}{m/s}$。音の速さ（約\qty{340}{m/s}）より速く、わたしたちは地面といっしょに動いています。赤道では緯線の半径が最も大きいので、約\qty{1680}{km/h}になります。

\begin{jibun}
地球を半径\qty{6400}{km}の球とします。（1）北緯45°、東経90°の地点の $(x,\ y,\ z)$ を求めましょう。（2）北緯40.8°の青森で、自転によって地面が動く速さは時速何kmですか。
\bkjans{（1）\nobreak $R\cos45°\approx4525$。$(0,\ 4525,\ 4525)$（km）。（2）\nobreak 半径 $6400\cos40.8°\approx4845$\,km、1周約\num{30441}\,km、24でわって時速約\qty{1268}{km}}
\end{jibun}

\begin{tsumazuki}{球座標の $\theta$ を、緯度と同じだと思う}
\bkng{北緯35.7°だから $\theta=35.7°$}
\bkok{$\theta$ は真上からはかるので $\theta=90°-35.7°=54.3°$}
緯度は赤道から、球座標の $\theta$ は北極（$z$ 軸）からはかります。はかり始める所がちがいます。
\end{tsumazuki}

\begin{reidai}[らせんの極方程式]
極方程式 $r=2\theta$（$\theta$ はラジアン）のらせんで、$\theta=\dfrac{\pi}{2}$、$\pi$、$2\pi$ のときの $r$ と $(x,\ y)$ を求めましょう。\par
\kangaekata $\theta$ を入れて $r$ を求め、①の式で $(x,\ y)$ にする。$\dfrac{\pi}{2}$ は90°、$\pi$ は180°、$2\pi$ は360°（4-2）。\par
\shiki $\theta=\dfrac{\pi}{2}$：$r=\pi$、\enspace $(\pi\cos90°,\ \pi\sin90°)=(0,\ \pi)$\par
\hspace*{2.6\zw}$\theta=\pi$：$r=2\pi$、\enspace $(2\pi\cos180°,\ 2\pi\sin180°)=(-2\pi,\ 0)$\par
\hspace*{2.6\zw}$\theta=2\pi$：$r=4\pi$、\enspace $(4\pi\cos360°,\ 4\pi\sin360°)=(4\pi,\ 0)$\par
\kotae $(0,\ \pi)$、$(-2\pi,\ 0)$、$(4\pi,\ 0)$（約 $(0,\ 3.14)$、$(-6.28,\ 0)$、$(12.57,\ 0)$）
\end{reidai}
1周した $\theta=2\pi$ の点は、はじめの向き（$x$ 軸の＋）にもどりますが、原点からの距離は $4\pi$ にふえています。だから、らせんになるのです。

\begin{jibun}
極方程式 $r=0.5\theta$ のらせんで、$\theta=\dfrac{\pi}{2}$、$\pi$、$2\pi$ のときの $(x,\ y)$ を求めましょう。
\bkjans{$(0,\ \dfrac{\pi}{4})$、$(-\dfrac{\pi}{2},\ 0)$、$(\pi,\ 0)$（約 $(0,\ 0.785)$、$(-1.57,\ 0)$、$(3.14,\ 0)$）}
\end{jibun}

\begin{tsumazuki}{極方程式に、度のまま入れる}
\bkng{$r=2\theta$ で、$\theta=90°$ のとき $r=180$}
\bkok{$\theta=\dfrac{\pi}{2}$ を入れて $r=\pi\approx3.14$}
$\theta$ が式の中で「数」として使われるときは、ラジアンで入れます（4-2）。電卓も RAD のモードにしましょう。
\end{tsumazuki}

\begin{column}{ロボットの手を、置きたい場所へ}
\noindent\begin{minipage}[c]{0.64\linewidth}
4-0の「逆運動学」の、いちばんかんたんな形を解いてみます。関節が1つで、うでの長さを伸び縮みさせられるロボットが、手を $(x,\ y)=(20,\ 15)$（cm）の部品の上に置きたいとします。うでの長さ $r$ と肩の角 $\theta$ は、$(x,\ y)$ を極座標に直せば求められます。
\bkivpline{$r=\sqrt{20^2+15^2}=\sqrt{625}=25$\qquad $\theta=\tan^{-1}\dfrac{15}{20}\approx36.9°$（右上なのでそのまま）}
関節が2つあるうでの角は、4-7の余弦定理で求めました。どちらも「置きたい場所から角を決める」計算です。
\end{minipage}\hfill
\begin{minipage}[c]{0.32\linewidth}\centering
\begin{tikzpicture}[x=1mm, y=1mm]
  \bkivpplane{-4}{32}{-4}{24}
  \draw[bkBlue, line width=2.6pt, line cap=round] (0,0) -- (24,18);
  \filldraw[fill=white, draw=bkBlue!70!black, line width=0.8pt] (0,0) circle (1.4);
  \fill[bkOrange] (24,18) circle (1.2);
  \draw[bkInk, line width=0.5pt] (7,0) arc (0:36.87:7);
  \node[font=\sffamily\footnotesize] at (18:10.5) {$\theta$};
  \node[font=\sffamily\footnotesize, text=bkBlue!80!black, anchor=south east, inner sep=1pt] at ($(12,9)+(-1,1)$) {$r$};
  \node[font=\sffamily\footnotesize, anchor=north, inner sep=1.5pt] at (24,16.4) {$(20,\ 15)$};
\end{tikzpicture}
\end{minipage}
\end{column}

\begin{reidai}[地球の表面にそった最短の距離]
地球を半径\qty{6400}{km}の球とします。A（北緯30°、東経0°）とB（北緯30°、東経90°）について、\par
（1）地球の中心から見た、AとBの向きのなす角を求めましょう。\par
（2）AからBまで、地球の表面にそった最短の距離（大円にそった距離）を求めましょう。\par
\kangaekata 中心からAの向きの単位ベクトルは、④の式で $R=1$ としたもの。なす角は内積で（4-13）。距離は弧の長さ＝半径×角（4-2。角はラジアン）。\par
\shiki（1）$\vec e_A=(\cos30°,\ 0,\ \sin30°)=(0.866,\ 0,\ 0.5)$\par
\hspace*{4.6\zw}$\vec e_B=(\cos30°\cos90°,\ \cos30°\sin90°,\ \sin30°)$\par
\hspace*{5.6\zw}$=(0,\ 0.866,\ 0.5)$\par
\hspace*{4.6\zw}なす角の cos は\enspace $\vec e_A\cdot\vec e_B=0+0+0.25=0.25$\par
\hspace*{4.6\zw}電卓（DEG）で $\cos^{-1}0.25\approx75.5°$\par
\hspace*{2.6\zw}（2）$75.5°$ をラジアンに直す\enspace $75.5\times\dfrac{\pi}{180}\approx1.318$\par
\hspace*{4.6\zw}$6400\times1.318\approx8436$\par
\kotae （1）約75.5°\quad（2）約\qty{8440}{km}
\end{reidai}
北緯30°の緯線にそって東へ進むと、緯線の半径 $6400\cos30°\approx5543$\,km の円の $\dfrac{1}{4}$ 周で、約\qty{8706}{km}です。大円にそうと、約\qty{270}{km}近いのです。大円は、北へふくらんだ道すじになります（コラム）。

\begin{jibun}
地球を半径\qty{6400}{km}の球とします。A（北緯60°、東経0°）とB（北緯60°、東経180°）の、中心から見たなす角と、大円にそった距離を求めましょう。
\bkjans{$\vec e_A=(0.5,\ 0,\ 0.866)$、$\vec e_B=(-0.5,\ 0,\ 0.866)$。内積 $-0.25+0.75=0.5$ で60°。$6400\times\dfrac{\pi}{3}\approx\qty{6702}{km}$（北極の上を通る）}
\end{jibun}

%% ---------------- 物理・工学 ----------------
\begin{butsuri}
\noindent\begin{minipage}[c]{0.62\linewidth}
\textbf{ロボット掃除機のセンサー}\par\noindent 本体の上で回るセンサー（LiDAR）は、光を出して、かべではね返ってくるまでの時間から、かべまでの距離 $r$ をはかります。センサーが1回転するあいだに、たとえば1°ごとに $(r,\ \theta)$ を360こはかり、$x=r\cos\theta$、$y=r\sin\theta$ で $(x,\ y)$ に直して、部屋の地図（右の図の点。図では15°ごと）を作ります。
\end{minipage}\hfill
\begin{minipage}[c]{0.35\linewidth}
\centering
\begin{tikzpicture}[x=1mm, y=1mm]
  \draw[bkGray!70, line width=1.2pt] (-18,-12) -- (22,-12) -- (22,14) -- (4,14) -- (4,20) -- (-18,20) -- cycle;
  \foreach \a/\tr in {0/22.00, 15/22.78, 30/25.40, 45/19.80, 60/16.17, 75/15.45, 90/20.00, 105/20.71, 120/23.09, 135/25.46, 150/20.78, 165/18.63, 180/18.00, 195/18.63, 210/20.78, 225/16.97, 240/13.86, 255/12.42, 270/12.00, 285/12.42, 300/13.86, 315/16.97, 330/24.00, 345/22.78} {
    \draw[bkRed!40, line width=0.3pt] (0,0) -- (\a:\tr);
    \fill[bkRed] (\a:\tr) circle (0.7);}
  \fill[bkGray!60] (0,0) circle (4);
  \fill[bkInk] (0,0) circle (1.5);
\end{tikzpicture}
\end{minipage}
\par\vspace{4pt}\noindent\textbf{原子の中の電子}\par\noindent 原子の中の電子が、どこにいやすいかを表す「電子の雲」の形は、原子核を原点とする球座標 $(r,\ \theta,\ \varphi)$ を使うと、とても簡単な式で書けます。周期表のしくみや、半導体の性質も、この形から説明されます（物理シリーズ 原-7）。
\end{butsuri}

%% ---------------- 章末問題 ----------------
\bkendproblems
\bkproblemgroup{基本}{3}

\bkproblem （ふく習）極座標が次の点の $(x,\ y)$ を求めましょう（4-8の例題2）。\quad（1）$(6,\ 60°)$\quad（2）$(4,\ 135°)$\quad（3）$(2,\ 300°)$
\bkpans{（1）\nobreak $(3,\ 3\sqrt3)$（約 $(3,\ 5.20)$）。（2）\nobreak $(-2\sqrt2,\ 2\sqrt2)$（約 $(-2.83,\ 2.83)$）。（3）\nobreak $(1,\ -\sqrt3)$（約 $(1,\ -1.73)$）}

\bkproblem 次の点の極座標を、$0°\leqq\theta<360°$ で求めましょう（\bkreidairef{1}）。\quad（1）$(-\sqrt3,\ 1)$\quad（2）$(1,\ -1)$\quad（3）$(-4,\ -3)$（$\theta$ は電卓で）
\bkpans{（1）\nobreak $r=2$、左上で $-30°+180°=150°$。$(2,\ 150°)$。（2）\nobreak $(\sqrt2,\ 315°)$。（3）\nobreak $r=5$、左下で $\tan^{-1}0.75\approx36.9°$、$+180°$ で $(5,\ 216.9°)$}

\bkproblem 半径\qty{0.4}{m}の円の上を、角速度\qty{5}{rad/s}で左回りに回る点が、いま $\theta=60°$ の所にいます。位置と速度の成分を求めましょう（\bkreidairef{2}）。
\bkpans{位置 $(0.2,\ 0.346)$。速さ2、$\vec e_\theta=(-0.866,\ 0.5)$。速度 $(-1.73,\ 1)$（m/s）}

\bkproblem （1）円柱座標 $r=8$、$\theta=240°$、$z=5$ の点の $(x,\ y,\ z)$ を求めましょう。（2）点 $(-4,\ -3,\ 10)$ の円柱座標を求めましょう（\bkreidairef{3}）。
\bkpans{（1）\nobreak 約 $(-4,\ -6.93,\ 5)$。（2）\nobreak $r=5$、左下で約216.9°。$(5,\ 216.9°,\ 10)$}

\bkproblem （1）球座標 $r=6$、$\theta=60°$、$\varphi=30°$ の点の $(x,\ y,\ z)$ を求めましょう。（2）点 $(2,\ 2,\ 2\sqrt2)$ の球座標を求めましょう（\bkreidairef{4}）。
\bkpans{（1）\nobreak $r\sin\theta\approx5.196$。約 $(4.5,\ 2.60,\ 3)$。（2）\nobreak $r=4$、$\cos\theta=\dfrac{\sqrt2}{2}$ で45°、$\varphi=45°$。$(4,\ 45°,\ 45°)$}

\bkproblem 極方程式 $r=\theta$ のらせんで、$\theta=\dfrac{\pi}{2}$、$\pi$、$2\pi$ のときの $(x,\ y)$ を求めましょう（\bkreidairef{6}）。
\bkpans{$(0,\ \dfrac{\pi}{2})$、$(-\pi,\ 0)$、$(2\pi,\ 0)$（約 $(0,\ 1.57)$、$(-3.14,\ 0)$、$(6.28,\ 0)$）}

\bkproblemgroup{チャレンジ}{5}

\bkproblem 地球を半径\qty{6400}{km}の球とします。（1）北緯60°、東経30°の地点の $(x,\ y,\ z)$ を求めましょう。（2）北緯26.2°の那覇で、自転によって地面が動く速さは時速何kmですか（\bkreidairef{5}）。
\bkpans{（1）\nobreak $R\cos60°=3200$。約 $(2771,\ 1600,\ 5543)$（km）。（2）\nobreak 半径 $6400\cos26.2°\approx5742$\,km。$2\pi\times5742\div24\approx$ 時速\qty{1503}{km}}

\bkproblem 地球を半径\qty{6400}{km}の球とします。A（北緯0°、東経0°）とB（北緯45°、東経90°）の、中心から見たなす角と、大円にそった距離を求めましょう（\bkreidairef{7}）。
\bkpans{$\vec e_A=(1,\ 0,\ 0)$、$\vec e_B=(0,\ 0.707,\ 0.707)$。内積0で90°。$6400\times\dfrac{\pi}{2}\approx\qty{10053}{km}$}

\bkproblem 球座標 $r=10$、$\theta=90°$、$\varphi=60°$ の点の $(x,\ y,\ z)$ を求めましょう。また、この点はどの平面の上にありますか（\bkreidairef{4}）。
\bkpans{$r\sin\theta=10$。約 $(5,\ 8.66,\ 0)$。$z=0$ なので $xy$ 平面の上}

\begin{butsuri}
\noindent\begin{minipage}[c]{0.7\linewidth}
\textbf{静止衛星}\par\noindent 天気の衛星「ひまわり」のような静止衛星は、赤道の上空約\qty{35800}{km}を、地球の自転と同じ24時間で1周します。地球の中心からの距離は $r=6400+\num{35800}=\num{42200}$（km）。地球の中心を原点にした極座標で見ると、$r$ は一定で、角 $\theta$ だけが24時間で360°ふえていく円運動です。速さは $v=r\omega$（4-8）なので、$v=\dfrac{2\pi\times\num{42200}}{24}\approx\qty{11000}{km/h}$。秒速約\qty{3.1}{km/s}。止まって見えるのに、じつはとても速く回っているのです。
\end{minipage}\hfill
\begin{minipage}[c]{0.27\linewidth}
\centering
\begin{tikzpicture}[x=0.85mm, y=0.85mm]
  \draw[bkGray, line width=0.5pt, dashed] (0,0) circle (19);
  \fill[bkBlue!15] (0,0) circle (6);
  \draw[bkBlue!70!black, line width=0.6pt] (0,0) circle (6);
  \draw[bkInk, line width=0.8pt] (0,0) -- (40:19);
  \fill[bkInk!70] (40:19) ++(-1.4,-1) rectangle ++(2.8,2);
  \draw[bk arrow, bkRed, line width=1.2pt] (40:19) -- ++(130:9);
  \node[font=\sffamily\footnotesize, anchor=north west, inner sep=1pt] at (20:10) {$r$};
  \node[font=\sffamily\footnotesize, text=bkRed!80!black, anchor=south, inner sep=2pt] at ($(40:19)+(130:9)$) {$v=r\omega$};
  \node[font=\sffamily\footnotesize, anchor=north] at (0,-6.5) {地球};
\end{tikzpicture}
\end{minipage}
\end{butsuri}

\begin{column}{地球は、少しつぶれた球}
この章では、地球を半径\qty{6400}{km}の球として計算しました。でも本当の地球は、赤道のまわりが少しふくらんだ形をしています。中心から赤道までは約\qty{6378}{km}、中心から北極までは約\qty{6357}{km}で、約\qty{21}{km}ちがいます。地球が1日に1回転しているので、赤道のあたりが外へふくらんだのです。GPSや地図アプリでは、このつぶれた形まで計算に入れて、緯度・経度・高さを $(x,\ y,\ z)$ に直しています。それでも、ちがいは半径の約 $\dfrac{1}{300}$ なので、この章の計算のように「球」と考えても、だいたいの距離は十分に求められます。
\begin{bkfigure}
\begin{tikzpicture}[x=1mm, y=1mm]
  \fill[bkBlue!12] (0,0) ellipse (22 and 16);
  \draw[bkBlue!70!black, line width=0.8pt] (0,0) ellipse (22 and 16);
  \draw[bkGray, line width=0.5pt, dashed] (0,0) circle (19);
  \draw[bkInk, line width=0.6pt, <->] (0,0) -- (22,0);
  \draw[bkInk, line width=0.6pt, <->] (0,0) -- (0,16);
  \node[font=\sffamily\footnotesize, anchor=north, inner sep=1.5pt] at (11,0) {\qty{6378}{km}};
  \node[font=\sffamily\footnotesize, anchor=west, inner sep=1.5pt] at (0,8) {\qty{6357}{km}};
  \node[font=\sffamily\footnotesize, anchor=south, inner sep=1.5pt] at (0,16.5) {北極};
  \node[font=\sffamily\footnotesize, anchor=west, inner sep=1.5pt] at (22.5,0) {赤道};
  \node[font=\sffamily\footnotesize, anchor=west, align=left] at (34,6) {つぶれ方を、わざと\\大きくしてかいた図。\\点線は、ふつうの球。};
\end{tikzpicture}
\end{bkfigure}
\end{column}

\begin{furikaeri}
\begin{checklist}
\item 極座標と $(x,\ y)$ を、区画に注意して行き来できる。\bkfurref{bk:4-15:kyoku}
\item $\vec e_r$、$\vec e_\theta$ を使って、円運動の速度を $r\omega\,\vec e_\theta$ と表せる。\bkfurref{bk:4-15:er}
\item 円柱座標・球座標で空間の点を表し、$(x,\ y,\ z)$ に直せる。\bkfurref{bk:4-15:enchu}
\end{checklist}
\end{furikaeri}

\bkfukidashi{delta}{同じ場所でも、言い方はいろいろ。形に合った座標を選べば、むずかしい問題も楽に解けるよ。}
